% !TeX root = ../Book1.tex
\begin{figure}[htbp]
	\centering
	\begin{tikzpicture}[x=3.65cm,y=3.65cm,>=Stealth,font=\small]
		\begin{scope}[shift={({0},0)}]
			\fill[blue!3] (0,0) rectangle (1,1);
			\draw[gray] (0,0) rectangle (1,1);
			\draw[blue!70!black,semithick] (0.0000,0.1400) -- (1.0000,0.6400);
			\draw[blue!70!black,semithick] (0.0000,0.6400) -- (0.7200,1.0000);
			\draw[blue!70!black,semithick] (0.7200,0.0000) -- (1.0000,0.1400);
			\fill[orange!85!black] (0,.14) circle[radius=.016];
			\node[below left] at (0,0) {$0$};
			\node[below] at (1,0) {$1$};
			\node[left] at (0,1) {$1$};
			\node at (.5,1.13) {有理斜率 $\alpha=1/2$};
			\node at (.5,-.18) {一周闭合轨迹};
		\end{scope}
		\begin{scope}[shift={({1.55},0)}]
			\fill[blue!3] (0,0) rectangle (1,1);
			\draw[gray] (0,0) rectangle (1,1);
			\draw[blue!70!black,semithick] (0.0000,0.1400) -- (0.6081,1.0000);
			\draw[blue!70!black,semithick] (0.6081,0.0000) -- (1.0000,0.5542);
			\draw[blue!70!black,semithick] (0.0000,0.5542) -- (0.3152,1.0000);
			\draw[blue!70!black,semithick] (0.3152,0.0000) -- (1.0000,0.9684);
			\draw[blue!70!black,semithick] (0.0000,0.9684) -- (0.0223,1.0000);
			\draw[blue!70!black,semithick] (0.0223,0.0000) -- (0.7294,1.0000);
			\draw[blue!70!black,semithick] (0.7294,0.0000) -- (1.0000,0.3826);
			\draw[blue!70!black,semithick] (0.0000,0.3826) -- (0.4365,1.0000);
			\draw[blue!70!black,semithick] (0.4365,0.0000) -- (1.0000,0.7969);
			\draw[blue!70!black,semithick] (0.0000,0.7969) -- (0.1436,1.0000);
			\draw[blue!70!black,semithick] (0.1436,0.0000) -- (0.8508,1.0000);
			\draw[blue!70!black,semithick] (0.8508,0.0000) -- (1.0000,0.2111);
			\draw[blue!70!black,semithick] (0.0000,0.2111) -- (0.5579,1.0000);
			\draw[blue!70!black,semithick] (0.5579,0.0000) -- (1.0000,0.6253);
			\draw[blue!70!black,semithick] (0.0000,0.6253) -- (0.2650,1.0000);
			\draw[blue!70!black,semithick] (0.2650,0.0000) -- (0.9721,1.0000);
			\draw[blue!70!black,semithick] (0.9721,0.0000) -- (1.0000,0.0395);
			\fill[orange!85!black] (0,.14) circle[radius=.016];
			\node[below left] at (0,0) {$0$};
			\node[below] at (1,0) {$1$};
			\node[left] at (0,1) {$1$};
			\node at (.5,1.13) {无理斜率 $\alpha=\sqrt2$};
			\node at (.5,-.18) {有限时间轨迹};
		\end{scope}
	\end{tikzpicture}
	\caption[环面上线性流的两种整体行为]{正方形的相对边分别识别，线段在右边离开后从左边进入，在上边离开后从下边进入。左图画出一次闭合的轨迹；右图只画出有限时间内的轨迹，稠密性由正文的整数逼近证明。两图中的每条局部线段都只有一维。}
	\label{g1:fig:torus-leaves}
\end{figure}

